% ---------------------------------------------------------------------------
% PDF/UA gate for the \ea ... \z front-end.
%
% The front-end opens lists and items in an order no other case produces:
% \ea opens an example, its list and its first item in one command, and \z
% closes a level whose \begingroup was opened by \lx@subopen@*.  That is
% exactly the shape in which a structure element gets opened at the wrong
% moment or left unclosed, and neither failure shows on the page:
% marked content straddling its parent fails veraPDF while passing every
% geometric assertion, and an element never closed is spec-valid -- veraPDF
% passes it -- while the rest of the document silently becomes its child.
% So this case is asserted on twice, by veraPDF and by struct_label_depths,
% which are complementary and neither of which is sufficient.
%
% The preamble is not shared: _preamble-ua.tex loads [lazy,gb4e] and
% already contains \begin{document}.  It mirrors it otherwise.
% "uncompress" lets the structure-tree helpers read the tree out of the raw
% bytes and does not affect validity; a dc:title is mandatory under
% ISO 14289-2 clause 8.11.1, hence the pdftitle.  Three LaTeX passes are
% needed for UA convergence under pdflatex and xelatex -- see PASSES.
%
% No footnote, for the reason ua.tex gives at length: on this TeX Live a
% footnote in a tagged UA-2 article emits untagged content with no linguexx
% loaded at all.  langsci-mixed.tex covers the footnote carve-out
% structurally, and examples/ua-demo end to end.
% ---------------------------------------------------------------------------
\DocumentMetadata{uncompress,lang=en,pdfversion=2.0,pdfstandard=ua-2,tagging=on}
\documentclass{article}
\usepackage{iftex}
\ifPDFTeX
  \usepackage[T1]{fontenc}
  \usepackage[utf8]{inputenc}
\else
  \usepackage{fontspec}
\fi
\usepackage[lazy,langsci]{linguexx}
\usepackage{hyperref}
\hypersetup{pdftitle={linguexx langsci front-end PDF/UA case},
            pdfdisplaydoctitle=true}
\pagestyle{empty}
\begin{document}
\section{UALSSECT the langsci front-end, tagged}

UALSINTRO ordinary running text precedes the first example.

\ea\label{uals:one} UALSMAIN a numbered example.
  \ea UALSALPHA a sub-example.
  \ex[*]{UALSBETA a judged sub-example.}
    \ea UALSROMAN a sub-sub-example.
    \z
  \z
\ex UALSTWO a second top-level example in the same block.
\z

\eal
  \ex UALSLISTA under a head with no text.
  \ex UALSLISTB the second letter.
\zl

\ea \gll UALSOBJ zwei Zeilen\\
        UALSGLOSS two lines\\
\glt UALSTRANS the free translation.
\z

\ex. UALSDOT an example in the dot syntax, in the same document.

%% The numbering variants, and the two constructs that had to be rewritten
%% out of math mode.  Both halves are asserted from the STRUCTURE TREE and
%% not from the page, because that is where each of them can go wrong
%% invisibly.
%%
%% A variant sets the printed label and the /ListNumbering class a screen
%% reader announces; a list labelled "A." whose class says LowerRoman is
%% well-formed PDF, passes veraPDF, and is simply false, so the page cannot
%% catch it and langsci-lists.tex asserts only the printed half.
%%
%% \exp and \atcenter are $'$ and $\vcenter{...}$ upstream.  Math here means
%% a Formula element, which is the PDF/UA-2 failure that forced the 0.12
%% removal of the old \altg -- so the tree must contain no Formula at all,
%% and a_langsci_ua asserts exactly that.  (The variants must sit in an exe
%% batch: they are environments, and the one-syntax-per-example rule keeps
%% environments out of an \ea example.)
\begin{exe}
\ex UALSVAR an example whose sub-levels carry the numbering variants.
 \begin{xlistA} \ex UALSUPALPHA an upper-alpha sub-example. \end{xlistA}
 \begin{xlistI} \ex UALSUPROMAN an upper-roman sub-example. \end{xlistI}
\exp{uals:one} UALSPRIME a primed cross-reference as a label.
\ex UALSCENTRE \atcenter{\hbox{UALSCA}\hbox{UALSCB}} centred beside it.
\end{exe}

%% The references are written \pLast / \pref, i.e. WITHOUT the
%% parentheses, and that is not cosmetic.  a_langsci_ua reads the nesting
%% depth of every top-level example number out of the structure tree, and
%% it finds them by their printed form: a "(5)" standing in running text is
%% indistinguishable from the fifth example's label and would be reported
%% as a number sitting at the wrong depth.  The paren-less forms take the
%% same anchor and make the same Link, so nothing about the tagging is
%% left uncovered -- and ua.tex asserts on the parenthesised spelling.
UALSREL a relative reference: \pLast, and a label: \pref{uals:one}.
\end{document}
